\section*{Capítulo \ref{ch:b2:reproductive-hormones} --- Controle hormonal da reprodução}

\begin{solution}{exo:b2:reproductive-hormones:1}
Hipotálamo (GnRH, pulsos) $\to$ adeno-hipófise (LH, FSH) $\to$
\omterm{def:b2:mammal-reproduction:testis}{testículo}: o LH sobre as \omterm{def:b2:mammal-reproduction:testis}{células de Leydig} (testosterona), o FSH sobre
as \omterm{def:b2:mammal-reproduction:testis}{células de Sertoli} (\omterm{def:b2:mammal-reproduction:testis}{espermatogênese}, inibina). Retroalimentação: a
testosterona inibe o GnRH e o LH; a inibina inibe o FSH.
\end{solution}

\begin{solution}{exo:b2:reproductive-hormones:2}
Ovário: \omterm{def:b2:reproductive-hormones:cycle}{fase folicular}, dias 1--14, estradiol do \omterm{def:b2:mammal-reproduction:ovary}{folículo} em
crescimento; \omterm{def:b2:mammal-reproduction:ovary}{ovulação}, dia 14, \omterm{def:b2:reproductive-hormones:cycle}{pico de LH}; \omterm{def:b2:reproductive-hormones:cycle}{fase lútea}, dias 15--28,
progesterona do \omterm{def:b2:mammal-reproduction:ovary}{corpo lúteo}. Útero: menstruação, dias 1--5 (queda dos
esteroides); fase proliferativa, dias 6--14 (estradiol); fase
secretora, dias 15--28 (progesterona).
\end{solution}

\begin{solution}{exo:b2:reproductive-hormones:3}
O estrógeno e o progestagênio diários mantêm o LH e o FSH baixos por
retroalimentação negativa, de modo que nenhum \omterm{def:b2:mammal-reproduction:ovary}{folículo} é recrutado,
nenhuma subida de estradiol ocorre e o pico nunca dispara; o
progestagênio também espessa o muco. Na semana sem pílula, os
esteroides caem e o \omterm{def:b2:reproductive-hormones:cycle}{endométrio}, construído sob a ação deles, descama
--- um sangramento de privação, e não uma menstruação após uma
\omterm{def:b2:mammal-reproduction:ovary}{ovulação}.
\end{solution}

\begin{solution}{exo:b2:reproductive-hormones:4}
O LH e o FSH sobem várias vezes (sem retroalimentação da testosterona
nem da inibina); os músculos, a libido e as glândulas acessórias
regridem. A testosterona restaura os caracteres e faz o LH baixar de
novo, mas não restaura a fertilidade: os \omterm{def:b2:mammal-reproduction:testis}{testículos} foram retirados e,
mesmo num homem com \omterm{def:b2:mammal-reproduction:testis}{testículos}, a testosterona exógena suprime o LH e,
portanto, a testosterona intratesticular de que a \omterm{def:b2:mammal-reproduction:testis}{espermatogênese}
precisa.
\end{solution}

\begin{solution}{exo:b2:reproductive-hormones:5}
$24/1.5 = 16$ pulsos por dia; $24/4 = 6$. Quando o \omterm{def:b2:mammal-reproduction:ovary}{corpo lúteo} morre,
os pulsos ainda são lentos, o que favorece o FSH: o FSH sobe primeiro
--- o hormônio certo, pois é o FSH que recruta a coorte seguinte de
\omterm{def:b2:mammal-reproduction:ovary}{folículos}.
\end{solution}

\begin{solution}{exo:b2:reproductive-hormones:6}
\qty{200}{pg/mL} $= \qty{200}{ng/L}$: $200\times 10^{-9}/272 =
\qty{7.4e-10}{mol/L} = \qty{740}{pmol/L}$. \qty{15}{ng/mL} $=
\qty{15}{\micro g/L}$: $15\times 10^{-6}/314 = \qty{48}{nmol/L}$. Num
microlitro: $7.4\times 10^{-16}$ mol $= 4.4\times 10^{8}$
moléculas.
\end{solution}

\begin{solution}{exo:b2:reproductive-hormones:7}
\omterm{def:b2:mammal-reproduction:ovary}{Ovulação} no dia $35 - 14 = 21$; período fértil do dia 18 ao dia 22.
Com ciclos irregulares, a \omterm{def:b2:reproductive-hormones:cycle}{fase folicular}, e portanto o dia da
\omterm{def:b2:mammal-reproduction:ovary}{ovulação}, não pode ser prevista a partir dos ciclos anteriores, e a
janela é perdida ou mal situada.
\end{solution}

\begin{solution}{exo:b2:reproductive-hormones:8}
Do dia 21 ao 26 são 2.5 duplicações: $2^{2.5} = 5.7$ UI/L --- o
bastante. Uma \omterm{def:b2:mammal-reproduction:implantation}{implantação} no dia 25 dá $1.4$ UI/L no dia 26; o \omterm{def:b2:mammal-reproduction:ovary}{corpo
lúteo}, já em regressão, não é resgatado, a progesterona cai e o
embrião se perde com a menstruação. A \omterm{def:b2:mammal-reproduction:implantation}{implantação} tardia é uma causa
comum de perda precoce.
\end{solution}

\begin{solution}{exo:b2:reproductive-hormones:9}
A hipófise responde ao padrão do GnRH, e não à sua quantidade: pulsos
horários sustentam o LH e o FSH, e a mesma dose diária infundida
continuamente os abole. O experimento descartou uma relação
dose--resposta simples. Um mês de agonista de ação prolongada: após
alguns dias de estímulo inicial, o LH cai a quase zero
(\omterm{prop:b2:reproductive-hormones:pulses}{dessensibilização dos receptores}), o estradiol a níveis de
pós-menopausa, o \omterm{def:b2:reproductive-hormones:cycle}{endométrio} se atrofia e nenhum ciclo ocorre --- uma
menopausa medicamentosa reversível.
\end{solution}

\begin{solution}{exo:b2:reproductive-hormones:10}
Sem hormônio: $R^* = 100/0.5 = 200$. Contínuo: $100/6.5 = 15$. Um
pulso de dois minutos remove $6\times 200\times(2/60) = 40$ receptores
(\qty{20}{\%}); na hora seguinte, o déficit diminui por um fator
$e^{-0.5}$, isto é, \qty{39}{\%} dele é recuperado. O estoque se
estabiliza onde a perda por pulso é igual à recuperação por hora:
cerca de 150 receptores antes de cada pulso, três quartos do máximo
--- dez vezes o nível contínuo.
\end{solution}

\begin{solution}{exo:b2:reproductive-hormones:11}
Menopausa: sem \omterm{def:b2:mammal-reproduction:ovary}{folículos}, portanto sem estradiol nem inibina, portanto
sem retroalimentação: o LH e o FSH ficam altos. Prolactinoma: a
prolactina desacelera o gerador de pulsos de GnRH, de modo que o LH é
baixo, os \omterm{def:b2:mammal-reproduction:ovary}{folículos} não são estimulados, o estradiol é baixo e não há
ciclo. Tumor de células da granulosa: o estradiol constantemente alto
suprime o LH por retroalimentação negativa (e, sem um \omterm{def:b2:mammal-reproduction:ovary}{folículo} em
maturação, não há nada para ovular), de modo que o ciclo para com
estradiol alto --- e o \omterm{def:b2:reproductive-hormones:cycle}{endométrio} cresce em excesso.
\end{solution}

\begin{solution}{exo:b2:reproductive-hormones:12}
O estradiol baixo ou breve reduz o GnRH e o LH; o estradiol alto e
sustentado (acima de \qty{200}{pg/mL} por mais de 36 horas) muda a
resposta do sistema kisspeptina--GnRH e da hipófise, que nesse meio
tempo acumulou LH, de modo que o mesmo hormônio passa a provocar uma
descarga. A \omterm{def:b2:mammal-reproduction:ovary}{ovulação} precisa ser de tudo ou nada e sincronizada: o
óvulo precisa ser liberado uma única vez, num único dia, quando um
\omterm{def:b2:mammal-reproduction:ovary}{folículo} está maduro. Um regulador gradual (LH proporcional ao
estradiol) daria uma luteinização gradual e parcial, \omterm{def:b2:mammal-reproduction:ovary}{folículos}
ovulando meio maduros ou vários de uma vez, e nenhum dia fértil
definido.
\end{solution}

\begin{solution}{pb:b2:reproductive-hormones:1}
\textbf{1.} Folicular, dias 1--14: estradiol subindo, progesterona
baixa. Lútea, dias 15--28: progesterona alta.
\textbf{2.} O pico atingiu o máximo no dia 14 (começou no dia 13); a
\omterm{def:b2:mammal-reproduction:ovary}{ovulação} ocorreu cerca de 36 horas após o início, no dia 15.
\textbf{3.} Dias 15 a 28: 14 dias, a vida do \omterm{def:b2:mammal-reproduction:ovary}{corpo lúteo}.
\textbf{4.} A progesterona eleva o valor de referência da
termorregulação em \qty{0.3}{\celsius} a \qty{0.4}{\celsius}; a subida
aparece um ou dois dias após a \omterm{def:b2:mammal-reproduction:ovary}{ovulação}, de modo que confirma que a
\omterm{def:b2:mammal-reproduction:ovary}{ovulação} ocorreu, mas não consegue anunciá-la.
\textbf{5.} Dias 11 a 14 (210, 260, 310, 220): três a quatro dias, mais
que as 36 horas que a chave exige.
\textbf{6.} Dias 12 a 16.
\textbf{7.} $310\times 10^{-9}/272 = \qty{1.14e-9}{mol/L} =
\qty{1140}{pmol/L}$; $16\times 10^{-6}/314 = \qty{51}{nmol/L}$.
\textbf{8.} $65/8 = 8$. Com uma meia-vida de uma hora, o LH cai à
metade a cada hora depois que a secreção para, de modo que um dia
depois ele está de volta perto do nível basal.
\textbf{9.} \omterm{def:b2:mammal-reproduction:ovary}{Ovulação} por volta do dia 22, ciclo de cerca de 36 dias.
\textbf{10.} \omterm{def:b2:mammal-reproduction:implantation}{Implantação} no dia 22; no dia 27 (cinco dias, 2.5
duplicações), cerca de \qty{5.7}{UI/L}.
\textbf{11.} A progesterona caindo a \qty{2}{ng/mL} no dia 27 significa
que o \omterm{def:b2:mammal-reproduction:ovary}{corpo lúteo} está regredindo no prazo: nada o resgatou.
\textbf{12.} Progesterona em torno de \qty{20}{ng/mL} e subindo,
estradiol em torno de \qty{200}{pg/mL}.
\textbf{13.} No fim do ciclo, o estradiol, a progesterona e a inibina
caem juntos e o FSH escapa de sua retroalimentação; à medida que a
coorte cresce, o estradiol e a inibina sobem e o FSH cai, e só o
\omterm{def:b2:mammal-reproduction:ovary}{folículo} com mais receptores de FSH (com mais células da granulosa)
continua crescendo com o FSH em queda --- os demais morrem.
\textbf{14.} $R^* = 200$ sem hormônio; $100/6.5 = 15$ sob hormônio
contínuo.
\textbf{15.} $6\times 200\times 2/60 = 40$ receptores, \qty{20}{\%};
constante de tempo $1/d = \qty{2}{h}$, de modo que \qty{39}{\%} do
déficit é recuperado em uma hora; o estoque se estabiliza em torno de
150 receptores, três quartos do máximo.
\textbf{16.} Os pulsos mantêm a hipófise com cerca de 150 receptores,
e ela responde a cada pulso com LH; uma infusão contínua leva o estoque
a 15 e as mesmas células silenciam, qualquer que seja a dose; os pulsos
restauram o estoque e a resposta.
\textbf{17.} No primeiro dia: o LH e a testosterona sobem (o agonista
estimula receptores que ainda estão lá --- é o estímulo inicial). Após
um mês: receptores dessensibilizados, LH perto de zero, testosterona em
níveis de castração, que é o objetivo terapêutico.
\textbf{18.} O FSH sobe (sem inibina), o LH fica normal; mais \omterm{def:b2:mammal-reproduction:ovary}{folículos}
sobrevivem à seleção a cada ciclo, com \omterm{def:b2:mammal-reproduction:ovary}{ovulações} múltiplas e gêmeos
dizigóticos.
\textbf{19.} Num ciclo natural, a queda do FSH mata todos os \omterm{def:b2:mammal-reproduction:ovary}{folículos}
exceto o dominante; um suprimento constante de FSH mantém a coorte
inteira crescendo até a maturidade.
\textbf{20.} Dez \omterm{def:b2:mammal-reproduction:ovary}{folículos} produzem muito mais que \qty{200}{pg/mL} de
estradiol precocemente, o que desencadearia um \omterm{def:b2:reproductive-hormones:cycle}{pico de LH} prematuro e
uma \omterm{def:b2:mammal-reproduction:ovary}{ovulação} antes da coleta; o antagonista bloqueia o GnRH e nenhum
pico pode ocorrer.
\textbf{21.} Na tabela, o pico atingiu o máximo no dia 14 e a \omterm{def:b2:mammal-reproduction:ovary}{ovulação}
ocorreu cerca de 36 horas após seu início; a hCG substitui o pico, e a
coleta às 35 horas apanha os ovócitos maduros, mas ainda dentro do
\omterm{def:b2:mammal-reproduction:ovary}{folículo}.
\textbf{22.} $10\times 0.7 = 7$ embriões, $7\times 0.4 = 2.8$
\omterm{def:b2:mammal-reproduction:implantation}{blastocistos}; $1 - 0.65^{2} = 0.58$.
\textbf{23.} O estrógeno e o progestagênio constantes mantêm o FSH e o
LH baixos: nenhum \omterm{def:b2:mammal-reproduction:ovary}{folículo} cresce, o estradiol nunca chega ao limiar,
não há pico, nem \omterm{def:b2:mammal-reproduction:ovary}{corpo lúteo}, nem subida da progesterona.
\textbf{24.} É preciso administrar testosterona para os caracteres
secundários e a libido; mas ela não restaura a alta concentração
intratesticular que as \omterm{def:b2:mammal-reproduction:testis}{células de Leydig} forneciam, de modo que a
\omterm{def:b2:mammal-reproduction:testis}{espermatogênese} continua suprimida --- que é justamente o objetivo.
\textbf{25.} \omterm{def:b2:mammal-reproduction:ovary}{Ovulação} no dia 15; \omterm{def:b2:reproductive-hormones:cycle}{fase lútea} de 14 dias; janela fértil
nos dias 12--16; pico desencadeado pelo estradiol acima de
\qty{200}{pg/mL} por mais de 36 horas, condição atendida nos dias
11--14.
\end{solution}
